\documentclass[14pt,a4paper]{extarticle}
\usepackage{fontspec}
\usepackage{polyglossia}
\setmainlanguage{russian}
\setmainfont{Times New Roman}
%\setmonofont{DejaVy Sans Mono}[Scale=MatchLowercase]
\usepackage{setspace}
\onehalfspacing
\usepackage{geometry}
\geometry{left=30mm, right=15mm, top=20mm, bottom=20mm}
\usepackage{indentfirst}
\setlength{\parindent}{1.25cm}
\setlength{\parskip}{0pt}
%\usepackage{url}
\usepackage{xurl}
\urlstyle{same} % для переносов в ссылках
\usepackage{hyperref}
\usepackage{booktabs}
\usepackage{ragged2e} % для justifying
\usepackage{wasysym} % для космических символов
\usepackage{graphicx}
\usepackage{float}
\usepackage{pdflscape}
\usepackage{listings}
\usepackage{xcolor}
\usepackage{array}

\lstset{
	language=C++,
	basicstyle=\small,
	keywordstyle=\color{blue},
	stringstyle=\color{red},
	commentstyle=\color{green!50!black},
	showstringspaces=false,
	breaklines=true,
	frame=single,
	numbers=left,
	numberstyle=\tiny\color{gray},
	stepnumber=1,
	keepspaces=true, % <-- важно для корректного отображения
	extendedchars=true,
	inputencoding=utf8
}

\begin{document}
	
	\sloppy
	\thispagestyle{empty}
	\centering{
		МИНИСТЕРСТВО НАУКИ И ВЫСШЕГО ОБРАЗОВАНИЯ\\
		РОССИЙСКОЙ ФЕДЕРАЦИИ\\
		Федеральное государственное автономное образовательное учреждение высшего образования <<Санкт-Петербургский политехнический университет Петра Великого>>\\[10pt]
		Институт компьютерных наук и кибербезопасности\\
		Высшая школа технологий искусственного интеллекта\\
		02.03.01 Математика и компьютерные науки\\
		\vspace{\fill}
		{\large
			Отчёт о лабораторной работе\\
			по дисциплине\\
			{\Large \bfseries Вычислительная математика}\\
			{\Large <<Вычисление приближённого значения табличной функции>>}\\
			Вариант 18\\
			\vspace{\fill}
		}
		
		{
			\flushright
			{\bfseries Выполнил:}\\
			студент группы 5130201/40003\\
			Джабраилов А.В.\\
			\hspace{9cm}Подпись: \hrulefill \\
			{\bfseries Принял:}\\
			К. физ.-мат. наук, доц., доц.\\
			Пак В.Г.\\
			\hspace{9cm}Подпись: \hrulefill \\
			
			\hspace{9cm}<<\rule{1.5cm}{0.4pt}>> \hrulefill\ 2026~г.
		}
		
		\centering{
			Санкт-Петербург, 2026
		}
	}
	
	\justifying
	
	\newpage
	\sloppy
	
	\section{Задание}
	Построить по таблице заданный интерполяционный многочлен, вычислить с помощью него приближённое значение функции в данной точке.
	
	\section{Цель работы}
	Научиться вычислять приближённые значения функции с помощью полиномиальной интерполяции и оценивать погрешность интерполяции.
	
	\section{Ход работы}
	Ниже приведена таблица функции для вычисления значения по многочлену Лагранжа (задание 1), по схеме Эйткена (задание 2) и многочлену Ньютона (задание 3):
	
	\begin{table}[h]
		\centering
		\begin{tabular}{|c|c|}
			\hline
			$x_i$ & $y_i$ \\ \hline
			-11,02 & 21,1230 \\ \hline
			-10,78 & 22,8279 \\ \hline
			-10,23 & 25,0046 \\ \hline
			-9,89 & 27,5928 \\ \hline
			-9,65 & 28,9933 \\ \hline
			-9,03 & 29,1933 \\ \hline
		\end{tabular}
	\end{table}
	
	Для выполнения заданий данной лабораторной работы была создана программа на языке C++ (см. раздел \hyperref[code]{"Приложение"}). 
	
	\newpage
	
	\subsection{Задание 1. По формуле Лагранжа}
	Найти приближённое значение функции при $x= -9,22$. Ниже приведён листинг кода функции, которая вычисляет значение по многочлену Лагранжа:
	
	\begin{lstlisting}[label=lag]
double lagrange(const std::vector<Point>& points, double x) {
	int n = points.size();
	double res = 0.0;
	for (int i = 0; i < n; ++i) {
		double xi = points[i].first;
		double yi = points[i].second;
		double t = 1.0;
// ...
	\end{lstlisting}
	
	Для данной таблицы значение в точке $x= -9,22$: $y(-9,22) = 29,18$.
	
	\subsection{Задание 2. По схеме Эйткена}
	Найти приближённое значение функции при $x= -9,22$. Ниже приведён листинг кода функции, которая вычисляет значение по схеме Эйткена:
	
	\begin{lstlisting}
double aitken(const std::vector<Point>& points, double x) {
	int n = points.size();
	std::vector<std::vector<double>> p(n, std::vector<double>(n, 0.0));
	for (int i = 0; i < n; ++i) {
		p[i][i] = points[i].second;
	}
// ...
	\end{lstlisting}
	
	Для данной таблицы значение в точке $x= -9,22$: $y(-9,22) = 29,18$.
	
	\subsection{Задание 3. По многочлену Ньютона с разделёнными разностями}
	Найти приближённое значение функции в точках $x_1= -10,44$ и $x_2= -9,41$. Ниже приведён листинг кода функции:
	
	\begin{lstlisting}
double newton1(const std::vector<Point>& points, double x) {
	int n = points.size();
	std::vector<Point> sorted_points = points;
	std::sort(sorted_points.begin(), sorted_points.end());
	std::vector<double> work(n);
	for (int i = 0; i < n; ++i) {
		work[i] = sorted_points[i].second;
	}
// ...
	\end{lstlisting}
	
	Для данной таблицы значения в указанных точках: $y(-10,44) = 23,83$; $y(-9,41) = 29,41$.
	
	\subsection{Задание 4. По многочлену Ньютона с конечными разностями}
	Для выполнения заданий №4 и №5 была дана следующая функция $f(x)$ с шагом $h$ от $a$ до $b$:
	
	\begin{table}[h]
		\centering
		\begin{tabular}{|c|c|c|c|}
			\hline
			$f(x)$ & $h$ & $a$ & $b$ \\ \hline
			$x^{5/\pi}$ & 0,4 & 0 & 2 \\ \hline
		\end{tabular}
	\end{table}
	
	Ниже приведён листинг кода функции для вычисления по многочлену Ньютона с конечными разностями:
	
	\begin{lstlisting}
double newton2(double x, double h, double a, double b) {
	auto f = [](double x) { return std::pow(x, 5.0 / M_PI); };
	int n = static_cast<int>((b - a) / h) + 1;
	std::vector<Point> points(n);
	for (int i = 0; i < n; ++i) {
		points[i] = {a + i * h, f(a + i * h)};
	}
	std::vector<std::vector<double>> diff;
	std::vector<double> y(n);
	for (int i = 0; i < n; ++i) {
		y[i] = points[i].second;
	}
// ...
	\end{lstlisting}
	
	Для данной таблицы значения в указанных точках: $y(-0,78) = 0,8578$; $y(1,43) = 1,77$.
	
	\subsection{Задание 5. По подходящему многочлену (Гаусса, Стирлинга или Бесселя)}
	У нас чётное число узлов (шаг 0.4, начало 0, конец 2), поэтому применяется многочлен Бесселя. Ниже приведён листинг года для вычисления ответа на данное задание:
	
	\begin{lstlisting}
double bessel(double x, double h, double a, double b) {
	auto f = [](double x) { return std::pow(x, 5.0 / M_PI); };
	
	int n = static_cast<int>((b - a) / h) + 1;
	std::vector<Point> points(n);
	for (int i = 0; i < n; ++i) {
		points[i] = {a + i * h, f(a + i * h)};
	}
// ...
	\end{lstlisting}
	
	Для данной таблицы значение в точке $x= 0,50$: $y(0,50) = 0,33$.
	
	\newpage
	
	\section{Результаты}
	Ответы:
	\begin{enumerate}
		\item $y(-9,22) = 29,18$
		\item $y(-9,22) = 29,18$
		\item $y(-10,44) = 23,83$; $y(-9,41) = 29,41$
		\item $y(-0,78) = 0,86$; $y(1,43) = 1,77$
		\item $y(0,50) = 0,33$
	\end{enumerate}
	
	\newpage
	
	\section{Приложение}
	\label{code}
	\begin{lstlisting}
#include <iostream>
#include <vector>
#include <algorithm>
#include <iomanip>
#include <cmath>

#define M_PI 3.14

using Point = std::pair<double, double>;

double lagrange(const std::vector<Point>& points, double x) {
	int n = points.size();
	double res = 0.0;
	
	for (int i = 0; i < n; ++i) {
		double xi = points[i].first;
		double yi = points[i].second;
		double t = 1.0;
		for (int j = 0; j < n; ++j) {
			if (j != i) {
				double xj = points[j].first;
				t *= (x - xj) / (xi - xj);
			}
		}
		res += yi * t;
	}
	
	return res;
}

double aitken(const std::vector<Point>& points, double x) {
	int n = points.size();
	std::vector<std::vector<double>> p(n, std::vector<double>(n, 0.0));
	
	for (int i = 0; i < n; ++i) {
		p[i][i] = points[i].second;
	}
	
	for (int k = 1; k < n; ++k) {
		for (int i = 0; i < n - k; ++i) {
			int j = i + k;
			double xi = points[i].first;
			double xj = points[j].first;
			p[i][j] = (p[i][j-1] * (xj - x) - p[i+1][j] * (xi - x)) / (xj - xi);
		}
	}
	
	return p[0][n-1];
}

double newton1(const std::vector<Point>& points, double x) {
	int n = points.size();
	std::vector<Point> sorted_points = points;
	std::sort(sorted_points.begin(), sorted_points.end());
	
	std::vector<double> work(n);
	for (int i = 0; i < n; ++i) {
		work[i] = sorted_points[i].second;
	}
	
	double res = work[0];
	double prod = 1.0;
	
	for (int j = 1; j < n; ++j) {
		for (int i = 0; i < n - j; ++i) {
			work[i] = (work[i+1] - work[i]) / (sorted_points[i+j].first - sorted_points[i].first);
		}
		prod *= (x - sorted_points[j-1].first);
		res += work[0] * prod;
	}
	
	return res;
}

double newton2(double x, double h, double a, double b) {
	auto f = [](double x) { return std::pow(x, 5.0 / M_PI); };
	
	int n = static_cast<int>((b - a) / h) + 1;
	std::vector<Point> points(n);
	for (int i = 0; i < n; ++i) {
		points[i] = {a + i * h, f(a + i * h)};
	}
	
	std::vector<std::vector<double>> diff;
	std::vector<double> y(n);
	for (int i = 0; i < n; ++i) {
		y[i] = points[i].second;
	}
	diff.push_back(y);
	
	for (int j = 1; j < n; ++j) {
		std::vector<double> row;
		for (int i = 0; i < n - j; ++i) {
			row.push_back(diff[j-1][i+1] - diff[j-1][i]);
		}
		diff.push_back(row);
	}
	
	double res;
	if (x > (a + b) / 2) {
		double t = (x - points.back().first) / h;
		res = diff[0].back();
		double term = 1;
		double fact = 1;
		for (int j = 1; j < n; ++j) {
			term *= (t + j - 1);
			fact *= j;
			res += diff[j].back() * term / fact;
		}
	} else {
		double t = (x - points[0].first) / h;
		res = diff[0][0];
		double term = 1;
		double fact = 1;
		for (int j = 1; j < n; ++j) {
			term *= (t - j + 1);
			fact *= j;
			res += diff[j][0] * term / fact;
		}
	}
	
	return res;
}

double bessel(double x, double h, double a, double b) {
	auto f = [](double x) { return std::pow(x, 5.0 / M_PI); };
	
	int n = static_cast<int>((b - a) / h) + 1;
	std::vector<Point> points(n);
	for (int i = 0; i < n; ++i) {
		points[i] = {a + i * h, f(a + i * h)};
	}
	
	std::vector<std::vector<double>> diff;
	std::vector<double> y(n);
	for (int i = 0; i < n; ++i) {
		y[i] = points[i].second;
	}
	diff.push_back(y);
	
	for (int j = 1; j < n; ++j) {
		std::vector<double> row;
		for (int i = 0; i < n - j; ++i) {
			row.push_back(diff[j-1][i+1] - diff[j-1][i]);
		}
		diff.push_back(row);
	}
	
	int k = 0;
	for (int i = 0; i < n - 1; ++i) {
		if (points[i].first <= x && x < points[i+1].first) {
			k = i;
			break;
		}
	}
	
	double x0 = points[k].first;
	double q = (x - x0) / h;
	
	double res = diff[0][k];
	
	if (n >= 2) {
		res += q * diff[1][k];
	}
	
	if (n >= 3 && k >= 1) {
		double coeff = q * (q - 1) / 2;
		double avg = (diff[2][k-1] + diff[2][k]) / 2;
		res += coeff * avg;
	}
	
	if (n >= 4 && k >= 1) {
		double coeff = q * (q - 1) * (q - 0.5) / 6;
		res += coeff * diff[3][k-1];
	}
	
	if (n >= 5 && k >= 2) {
		double coeff = q * (q * q - 1) * (q - 2) / 24;
		double avg = (diff[4][k-2] + diff[4][k-1]) / 2;
		res += coeff * avg;
	}
	
	if (n >= 6 && k >= 2) {
		double coeff = q * (q * q - 1) * (q - 0.5) * (q - 2) / 120;
		res += coeff * diff[5][k-2];
	}
	
	return res;
}

int main() {
	std::vector<Point> points = {
		{-11.02, 21.1230},
		{-10.78, 22.8279},
		{-10.23, 25.0046},
		{-9.89, 27.5928},
		{-9.65, 28.9933},
		{-9.03, 29.1933},
	};
	double h = 0.4;
	double a = 0;
	double b = 2;
	
	double x = -9.22;
	std::cout << " " << x << ": " << std::fixed << std::setprecision(2) << lagrange(points, x) << std::endl;
	std::cout << " " << x << ": " << std::fixed << std::setprecision(2) << aitken(points, x) << std::endl;
	std::cout << std::endl;
	
	double x1 = -10.44;
	double x2 = -9.41;
	std::cout << " " << x1 << ": " << std::fixed << std::setprecision(2) << newton1(points, x1) << std::endl;
	std::cout << " " << x2 << ": " << std::fixed << std::setprecision(2) << newton1(points, x2) << std::endl;
	std::cout << std::endl;
	
	x1 = -0.78;
	x2 = 1.43;
	x = 0.50;
	std::cout << " " << x1 << ": " << std::fixed << std::setprecision(2) << newton2(x1, h, a, b) << std::endl;
	std::cout << " " << x2 << ": " << std::fixed << std::setprecision(2) << newton2(x2, h, a, b) << std::endl;
	std::cout << " " << x << ": " << std::fixed << std::setprecision(2) << bessel(x, h, a, b) << std::endl;
	std::cout << std::endl;
	
	return 0;
}
	\end{lstlisting}
	
\end{document}